\begin{tikzpicture}[
  every node/.style = {font=\small},
  idx/.style = {
    draw, rounded corners=3pt,
    minimum width=0.85cm, minimum height=0.60cm,
    inner sep=3pt, align=center
  },
  kept/.style  = {idx, fill=blue!12,  draw=blue!50!black},
  drop/.style  = {idx, fill=red!12,   draw=red!60!black, text=red!70!black},
  lbl/.style   = {font=\footnotesize\itshape, text=black!55},
  arr/.style   = {-{Stealth[length=6pt]}, thick, black!60},
  grpbox/.style = {
    draw=orange!70!black, rounded corners=5pt,
    fill=orange!6, inner sep=10pt, dashed, thick
  },
  derivbox/.style = {
    draw=teal!60!black, rounded corners=5pt,
    fill=teal!8, inner sep=8pt, thick
  },
]

%% ---------------------------------------------------------------
%% Step 0: root variable label at top
%% ---------------------------------------------------------------
\node[font=\normalsize\bfseries] (title) at (2.5, 0.5)
  {Index suppression};

%% ---------------------------------------------------------------
%% Root variable row (shown once, at top)
%% ---------------------------------------------------------------
\node[lbl] (rlbl) at (2.5, -0.2) {root variable $\phi$};

\node[font=\normalsize]          at (0.0, -1.0) {$\phi$};
\node[kept] (ra) at (0.8,  -1.0) {$i_1$};
\node[kept] (rb) at (1.75, -1.0) {$\cdots$};
\node[kept] (rc) at (2.7,  -1.0) {$i_k$};
\node[drop] (rd) at (3.7,  -1.0) {$j_1$};
\node[drop] (re) at (4.65, -1.0) {$\cdots$};
\node[drop] (rf) at (5.6,  -1.0) {$j_l$};

% braces
\draw[decorate, decoration={brace, amplitude=4pt},
      blue!50!black, thick]
  ($(ra.north west)+(0,0.06)$) -- ($(rc.north east)+(0,0.06)$)
  node[midway, above=5pt, font=\footnotesize, text=blue!60!black]
    {retained};

\draw[decorate, decoration={brace, amplitude=4pt},
      red!60!black, thick]
  ($(rd.north west)+(0,0.06)$) -- ($(rf.north east)+(0,0.06)$)
  node[midway, above=5pt, font=\footnotesize, text=red!60!black]
    {suppressed};

%% ---------------------------------------------------------------
%% Arrow down: "for each fixed (i1,...,ik)"
%% ---------------------------------------------------------------
\draw[arr] (2.8,-1.45) -- (2.8,-2.15)
  node[midway, right=5pt, lbl] {fix retained values $(i_1,\ldots,i_k)$};

%% ---------------------------------------------------------------
%% Aggregation group G: list of root instances varying j's
%% ---------------------------------------------------------------

% Four representative instances inside the group box
\node[font=\small] (g1) at (1.4, -2.9)
  {$\phi_{i_1,\ldots,i_k,\;j_1^{(1)},\ldots,j_l^{(1)}}$};
\node[font=\small] (g2) at (1.4, -3.55)
  {$\phi_{i_1,\ldots,i_k,\;j_1^{(2)},\ldots,j_l^{(2)}}$};
\node[font=\small] (g3) at (1.4, -4.2)
  {$\vdots$};
\node[font=\small] (g4) at (1.4, -4.7)
  {$\phi_{i_1,\ldots,i_k,\;j_1^{(s)},\ldots,j_l^{(s)}}$};

% Fit a box around them
\begin{pgfonlayer}{background}
  \node[grpbox, fit=(g1)(g2)(g3)(g4)] (gbox) {};
\end{pgfonlayer}

% Label G below the box
\node[font=\normalsize, text=orange!80!black, below=4pt of gbox]
  {$\mathcal{G}_{i_1,\ldots,i_k}$};

%% ---------------------------------------------------------------
%% Arrow right: aggregation rule
%% ---------------------------------------------------------------
\draw[arr] (gbox.east) -- ++(2.2, 0)
  node[midway, above=4pt, lbl] {aggregation rule}
  node[right] (deriv-anchor) {};

%% ---------------------------------------------------------------
%% Derived variable box
%% ---------------------------------------------------------------
\node[derivbox, right=2.2cm of gbox] (dbox)
  {$\phi_{i_1,\ldots,i_k}$};

\node[font=\footnotesize\itshape, text=teal!60!black, below=4pt of dbox]
  {derived variable};

%% ---------------------------------------------------------------
%% Annotation: "all j-tuples in J1 x ... x Jl"
%% ---------------------------------------------------------------
\node[lbl, below=0.55cm of gbox, yshift=-0.15cm]
  {$j_n^{(\cdot)} \in J_n,\quad n=1,\ldots,l$};

\end{tikzpicture}
